\documentclass[UTF8,11pt]{ctexart}

\usepackage[a4paper,margin=2.4cm]{geometry}
\usepackage{amsmath,amssymb}
\usepackage{xcolor}
\usepackage{enumitem}
\usepackage{fancyhdr}
\usepackage{hyperref}

\definecolor{courseblue}{RGB}{43,49,139}
\hypersetup{colorlinks=true,linkcolor=courseblue,urlcolor=courseblue}
\setlength{\parindent}{2em}
\setlength{\parskip}{0.35em}
\setlist{itemsep=0.25em,topsep=0.35em}

\pagestyle{fancy}
\setlength{\headheight}{14pt}
\fancyhf{}
\lhead{复变函数补充阅读}
\rhead{03}
\cfoot{\thepage}
\renewcommand{\headrulewidth}{0.4pt}

\newcommand{\R}{\mathbb{R}}
\newcommand{\dif}{\mathop{}\!\mathrm{d}}

\title{\textcolor{courseblue}{补充阅读 03：第二类曲线积分复习}}
\author{}
\date{}

\begin{document}
\maketitle
\vspace{-2.2em}

\section{为什么复积分会用到第二类曲线积分}

设
\[
f(z)=u(x,y)+iv(x,y),\qquad z=x+iy.
\]
由于
\[
\dif z=\dif x+i\dif y,
\]
所以
\[
\begin{aligned}
f(z)\dif z
&=(u+iv)(\dif x+i\dif y)\\
&=(u\dif x-v\dif y)+i(v\dif x+u\dif y).
\end{aligned}
\]
因此，复积分可以拆成两个实的第二类曲线积分：
\[
\boxed{
\int_C f(z)\dif z
=\int_C(u\dif x-v\dif y)
+i\int_C(v\dif x+u\dif y).
}
\]
理解第二类曲线积分的参数化、方向，以及它在什么条件下与路径无关，是学习复积分的直接准备。

\section{定义与直观意义}

设有向曲线 $C$ 由参数方程
\[
x=x(t),\qquad y=y(t),\qquad \alpha\leq t\leq\beta
\]
给出，其中 $x(t),y(t)$ 分段光滑。第二类曲线积分定义为
\[
\boxed{
\int_C P(x,y)\dif x+Q(x,y)\dif y
=\int_\alpha^\beta
\big[P(x(t),y(t))x'(t)+Q(x(t),y(t))y'(t)\big]\dif t.
}
\]

若把
\[
\mathbf F=(P,Q),\qquad \dif\mathbf r=(\dif x,\dif y)
\]
分别看作平面向量场和沿曲线的微小位移，那么
\[
P\dif x+Q\dif y=\mathbf F\cdot\dif\mathbf r.
\]
因此，第二类曲线积分可以理解为向量场沿路径所做功的累积。

它与曲线的方向有关。若 $C^{-}$ 表示把 $C$ 反向，则
\[
\int_{C^-}P\dif x+Q\dif y
=-\int_C P\dif x+Q\dif y.
\]

\section{参数化计算的四个步骤}

计算
\(
\int_C P\dif x+Q\dif y
\)
时，可以依次进行：
\begin{enumerate}
  \item 按题目规定的方向写出参数方程
        \[
        x=x(t),\qquad y=y(t),\qquad \alpha\leq t\leq\beta.
        \]
  \item 计算微分
        \[
        \dif x=x'(t)\dif t,
        \qquad
        \dif y=y'(t)\dif t.
        \]
  \item 把 $x,y,\dif x,\dif y$ 全部代入被积表达式。
  \item 按参数区间计算普通定积分。
\end{enumerate}

这里不能只代入 $x(t),y(t)$ 而忘记替换 $\dif x,\dif y$。参数区间的顺序必须与曲线方向一致。

\section{例 1：沿线段计算}

设 $C$ 是从 $(0,0)$ 到 $(1,2)$ 的线段，计算
\[
I=\int_C(x+y)\dif x+(x-y)\dif y.
\]
取参数方程
\[
x=t,\qquad y=2t,\qquad 0\leq t\leq1.
\]
于是
\[
\dif x=\dif t,\qquad \dif y=2\dif t,
\]
从而
\[
\begin{aligned}
I
&=\int_0^1\big[(t+2t)\cdot1+(t-2t)\cdot2\big]\dif t\\
&=\int_0^1 t\dif t=\frac12.
\end{aligned}
\]
如果把线段方向反过来，积分值变为 $-1/2$。

\section{例 2：沿圆周计算}

设 $C$ 是单位圆 $x^2+y^2=1$，方向为逆时针，计算
\[
I=\oint_C-y\dif x+x\dif y.
\]
取
\[
x=\cos t,\qquad y=\sin t,\qquad 0\leq t\leq2\pi.
\]
则
\[
\dif x=-\sin t\dif t,\qquad
\dif y=\cos t\dif t,
\]
所以
\[
\begin{aligned}
I
&=\int_0^{2\pi}
\big[(-\sin t)(-\sin t)+(\cos t)(\cos t)\big]\dif t\\
&=\int_0^{2\pi}1\dif t=2\pi.
\end{aligned}
\]
若改为顺时针方向，积分值为 $-2\pi$。

\section{全微分与路径无关}

如果存在函数 $\Phi(x,y)$，使得
\[
P=\Phi_x,
\qquad
Q=\Phi_y,
\]
那么
\[
P\dif x+Q\dif y=\dif\Phi
\]
称为一个全微分。此时，对于从点 $A$ 到点 $B$ 的任意分段光滑曲线 $C$，
\[
\boxed{
\int_C P\dif x+Q\dif y=\Phi(B)-\Phi(A).
}
\]
积分只由起点和终点决定，与具体路径无关。特别地，对任意闭曲线，
\[
\oint_C P\dif x+Q\dif y=0.
\]

若 $P,Q$ 的一阶偏导数连续，并且定义域单连通，则
\[
P_y=Q_x
\]
是 $P\dif x+Q\dif y$ 为全微分的常用判据。

第三章例 1.1 中出现的两个表达式就是全微分：
\[
x\dif x-y\dif y
=\dif\left(\frac{x^2-y^2}{2}\right),
\]
\[
y\dif x+x\dif y=\dif(xy).
\]
因此，相应的积分都只取决于端点。

\section{格林公式}

设 $C$ 是分段光滑的简单闭曲线，并按正方向围成区域 $D$。若 $P,Q$ 具有连续的一阶偏导数，则格林公式为
\[
\boxed{
\oint_C P\dif x+Q\dif y
=\iint_D(Q_x-P_y)\dif x\dif y.
}
\]
这里“正方向”是指沿 $C$ 前进时，区域 $D$ 始终位于左侧。对单个外边界，它就是逆时针方向。

当 $P_y=Q_x$ 时，右侧为零，于是闭路积分为零。这正是第三章用格林公式解释柯西定理合理性的基础。

\section{与复积分的直接连接}

设 $C$ 由
\[
z=z(t)=x(t)+iy(t),
\qquad \alpha\leq t\leq\beta
\]
给出，则
\[
\dif z=z'(t)\dif t.
\]
将两个第二类曲线积分合并，就得到复积分的参数公式：
\[
\boxed{
\int_C f(z)\dif z
=\int_\alpha^\beta f(z(t))z'(t)\dif t.
}
\]

例如，在单位圆 $C: z=e^{it}$ $(0\leq t\leq2\pi)$ 上，
\[
\oint_C\overline z\dif z
=\int_0^{2\pi}e^{-it}\,ie^{it}\dif t
=2\pi i.
\]
从第二类曲线积分的角度看，$\overline z=x-iy$，因此
\[
\oint_C\overline z\dif z
=\oint_C(x\dif x+y\dif y)
+i\oint_C(-y\dif x+x\dif y).
\]
第一个积分是闭路上的全微分积分，等于 $0$；第二个积分由例 2 得到 $2\pi$，所以结果仍为 $2\pi i$。

\section{常见错误}

\begin{enumerate}
  \item \textbf{忽略方向。} 同一条曲线反向后，第二类曲线积分变号。
  \item \textbf{忘记替换微分。} 代入 $x(t),y(t)$ 后，还必须代入 $\dif x=x'(t)\dif t$、$\dif y=y'(t)\dif t$。
  \item \textbf{把第二类曲线积分与第一类曲线积分混淆。} $P\dif x+Q\dif y$ 与方向有关，而含 $\dif s$ 的第一类曲线积分通常与方向无关。
  \item \textbf{认为闭曲线积分必为零。} 只有在满足全微分、格林公式中的相应条件或其他专门定理时，闭路积分才一定为零。
  \item \textbf{未经检查就使用端点代入。} 只有确认 $P\dif x+Q\dif y=\dif\Phi$ 后，才能直接计算 $\Phi(B)-\Phi(A)$。
\end{enumerate}

\section{自测}

\begin{enumerate}
  \item 设 $C$ 是从 $(0,0)$ 到 $(2,1)$ 的线段，计算
        \[
        \int_C y\dif x+x\dif y.
        \]
  \item 设 $C$ 是半径为 $R$ 的圆，方向为逆时针，计算
        \[
        \oint_C -y\dif x+x\dif y.
        \]
  \item 判断
        \[
        (2x+y)\dif x+(x+2y)\dif y
        \]
        是否为全微分，并计算它沿任意一条从 $(0,0)$ 到 $(1,1)$ 的曲线的积分。
  \item 为什么仅仅知道 $C$ 是闭曲线，不能断定
        \(
        \oint_C P\dif x+Q\dif y=0
        \)？
\end{enumerate}

\subsection*{参考答案}

\begin{enumerate}
  \item 因为 $y\dif x+x\dif y=\dif(xy)$，所以积分为 $2$。
  \item 参数化 $x=R\cos t$、$y=R\sin t$，得到 $2\pi R^2$。
  \item 它等于
        \[
        \dif(x^2+xy+y^2),
        \]
        因此积分为 $3$。
  \item 一般的第二类曲线积分与路径有关。只有在被积表达式为全微分，或满足格林公式、柯西定理等相应条件时，闭路积分才必为零。
\end{enumerate}

\end{document}
